% =====================================================================================
%  Lecture 3  —  21 September 2026
% =====================================================================================
\chapter{Why do big models generalize?}\label{ch:L3}
\setlecturevideo{q_Dj4-QNJV8}
\lectureinfo{21 September 2026}{1 h 46 min}{q_Dj4-QNJV8}{scaling revisited, generalization, compression, singular learning theory}

\begin{glance}
\begin{itemize}
  \item A list of \emph{empirical} puzzles in deep learning that lack a theory: scaling laws, generalization, grokking, shared representations, in-context learning.
  \item The spectral model of Lecture~2 revisited: what it assumes, and a test against data (image classification on CIFAR; still open for language models).
  \item A heuristic for scaling with \emph{model size}: the dimension of the data manifold gives $L\propto N^{-4/d}$.
  \item \textbf{Generalization:} the classical bias--variance picture, and Zhang et al.'s experiment that breaks it.
  \item \textbf{Learning as compression:} Bayesian prediction, log-loss, Kraft and Shannon, and Solomonoff's prior, under which regret is bounded by the Kolmogorov complexity of the truth; then Rissanen's $\tfrac d2\log n$ and Watanabe's singular learning theory.
\end{itemize}
\end{glance}

\begin{caution}[About the source for this lecture]
YouTube provides no English captions for this recording (the auto-dubbed tracks are in other languages). I worked from machine translations of the captions into English via four different languages, cross-checked against each other, and reconstructed the mathematics from them. Technical terms were occasionally mistranslated (for example, ``in-context learning'' came out as ``contextual learning''); I have normalized them. The names of authors and papers are the ones I could identify with confidence; those I am less sure of are marked\flag.
\end{caution}

\begin{figure}[tb]
\centering
\mapbar{6363}{900}{%
  \mapseg{0}{688}{slate}\mapseg{688}{2081}{teal}\mapseg{2081}{2720}{accent}\mapseg{2720}{3168}{amber}
  \mapseg{3168}{4045}{crimson}\mapseg{4045}{5546}{plum}\mapseg{5546}{6363}{moss}
  \mapin{344}{3.1}\mapin{1384}{3.2 \ Spectral law revisited}\mapin{2400}{3.3}\mapin{2944}{3.4}\mapin{3606}{3.5 \ Classical view}\mapin{4795}{3.6 \ Compression}\mapin{5954}{3.7 \ Parametric}
  \maplabel{344}{Puzzles}{3.2}{south west}
  \maplabel{2400}{Data check}{7.2}{south west}
  \maplabel{2944}{$N^{-4/d}$}{3.2}{south west}
}
\caption{\textbf{Lecture map} (1\,h\,46\,min, to scale). The first part (3.1--3.4) is about scaling laws, the second (3.5--3.7) about generalization.}
\end{figure}

%% ------------------------------------------------------------------------------------
\section{Observations without a theory}\label{sec:L3-observations}
\tsm{0:01}
After some announcements (a Monday seminar on \emph{corrigibility}, i.e.\ how to train a system persistent at its goals so that the goals can still be updated; a guest lecture by Anson Ho of Epoch AI on Wednesday; everything is recorded) the instructor sets the plan.
The course outline for this week was \emph{theories of learning, deep learning in particular}, which was ``probably a bit ambitious'' for one week. The coverage will be incomplete, and will not match the reading material exactly; the field is large, and it makes sense to focus on what one finds interesting and understand it as well as possible rather than everything.
Being new to the area himself, the instructor warns that his picks may not be the aspects most relevant for safety.

\paragraph{The analogy again.}\tsm{4:14}
We are in a situation like that of physics in the early twentieth century: interesting observations (the photoelectric effect, Michelson--Morley) but no complete explanation. \emph{We have no good theory with which to reason about thinking machines}, and progress on such a theory would be valuable for understanding risks.
What are the empirical observations?

\begin{table}[tb]
\centering\small
\renewcommand{\arraystretch}{1.28}
\begin{tabular}{@{}>{\raggedright\arraybackslash}p{26mm}>{\raggedright\arraybackslash}p{41mm}>{\raggedright\arraybackslash}p{31mm}>{\raggedright\arraybackslash}p{14mm}@{}}
\toprule
\textbf{Phenomenon} & \textbf{What is observed} & \textbf{Why safety cares} & \textbf{Here?}\\
\midrule
Scaling laws & Loss is a power law in compute, data, parameters (Lecture~2). & Forecasting capabilities; compute governance. & recap, \S\ref{sec:L3-recap}--\ref{sec:L3-params}\\
Generalization & A network trained on some data does very well on data it has never seen; the classical theory had expected this to be hard. & What can and cannot be learned; behavior off the training distribution. & main topic, \S\ref{sec:L3-classical}--\ref{sec:L3-singular}\\
Grokking & Perfect training accuracy early, poor test accuracy for a long time, then an abrupt jump, like a phase transition. & Sudden capability jumps. & later\\
Platonic representations & Networks with different architectures trained on different data have surprisingly similar internal features. & Transferability of interpretability tools. & later\\
Linear representations & Features are largely linear directions in activation space. & Probes, steering, monitoring. & later\\
Entangled concepts (``subliminal learning'') & Many concepts are intertwined in ways that let traits pass through apparently unrelated data\flag. & Hidden transmission of behavior. & later\\
In-context learning & Given examples in the prompt, a model infers the pattern without any weight update; absent in small or early models, then it appears suddenly. & Learning \emph{at inference}; jumps. & later\\
\bottomrule
\end{tabular}
\caption{The empirical puzzles listed at the start of the lecture, with the instructor's remark that for each of them ``we have some partial theories, but not a complete understanding''. Only scaling and generalization are treated today. Pointers\dagmark: grokking \citep{power2022}; Platonic representations \citep{huh2024}; linear representations \citep{park2024}; subliminal learning \citep{cloud2025}; in-context learning \citep{olsson2022}.}
\end{table}

\paragraph{Generalization is the fundamental observation.}\tsm{5:14}
The most basic fact about deep learning is that a network trained on some data works very well on data it has not seen. This runs against the intuition of earlier theory, which supposed it hard to make systems generalize and found, surprisingly, that it was easy. It is \emph{not quite right to say that it contradicted the theory},
because the theory focused on the wrong situation: it was overly pessimistic about what can and cannot be learned.

\paragraph{Grokking.}\tsm{6:19}
In the instructor's description: when a model is trained on certain tasks it may reach perfect accuracy on the training set but not on the test set; the loss nevertheless continues to decrease (loss is not accuracy); and then, suddenly, performance on unseen data improves dramatically. Such a phase after memorization followed by a sudden change can be described as a \emph{phase transition},
and some physics-like theories explain the resemblance. From the safety viewpoint it matters to understand \emph{when a network suddenly becomes much better at a task than before}\,\citep{power2022}.

\begin{discussion}[Where does the word ``grokking'' come from?]
The intuition: when learning a new language one can sit down and study the grammar, but that is not how humans actually learn; they hear and memorize, and suddenly everything falls into place and they understand the underlying rule. The word describes that moment. A student adds that it spread from a hacker subculture, where it meant to understand deeply.\dagmark{} (The word itself comes from Heinlein's novel \emph{Stranger in a Strange Land}, 1961.)
\end{discussion}

%% ------------------------------------------------------------------------------------
\section{The spectral law revisited}\label{sec:L3-recap}
\tsm{11:28}
\paragraph{Recap.}\tsm{13:01}
Scaling laws were observed over very large ranges, with Chinchilla looking at a wider range than Kaplan; the loss falls linearly on a log scale as compute, data or parameters grow, and once compute is fixed there is a trade-off between data and parameters, $C\propto N\cdot D$. The instructor repeats the solvable model of Lecture~2, in one sentence each:
\begin{itemize}
  \item the parameters $\theta\in\R^n$ define a function $f_\theta$ on the data space; one wants $f_\theta\to f^*$ by following the gradient of the loss; the true loss would be over \emph{all} data but one uses the empirical loss, and in the scaling regime the training and test losses are so similar as to be indistinguishable (the reason is the subject of the second half of the lecture);
  \item the \emph{neural tangent kernel} $K_\theta(x,x')=\ip{\nabla_\theta f_\theta(x)}{\nabla_\theta f_\theta(x')}$ measures how similar the gradient sensitivities of two inputs are, and gradient flow becomes $\partial_tg=-T\,g$ with $T$ self-adjoint, positive and compact (Lemma~\ref{lem:T});
  \item expanding the target in the eigenfunctions of $T$ gives an ODE whose solution converges, mode by mode, at a rate set by the eigenvalue.
\end{itemize}
This time the exponents are named differently: eigenvalues $\lambda_k\sim k^{-\alpha}$ and target coefficients $a_k^2\sim k^{-1-s}$ with $s>0$ (so that $f^*$ has finite $L^2$ norm). The result is
\[
   L(t)\;\asymp\;t^{-s/\alpha},\qquad\text{i.e. Proposition~\ref{prop:powerlaw} with } s=\alpha\beta .
\]
The quantity $t^{1/\alpha}$ is the transition index between the eigenfunctions that have been learned and those that have not, and $e^{-2\lambda_kt}$ is the factor that separates them. The instructor re-derives the bounds along the lines of Section~\ref{sec:L2-power}: the tail $k\ge m$ gives the lower bound, the whole sum the matching upper bound. (He stumbles over which exponent is which, a mix-up that he catches and corrects: $\alpha$ is the decay of the eigenvalues $\lambda_k$, $s$ that of the target coefficients $a_k^2$. If the target is concentrated on the first few eigenfunctions, learning is fast.)

\begin{discussion}[Questions from the class]
\emph{What architecture does this need?} None: nothing was assumed about the architecture except that the network evolves by gradient flow; it could be as simple as a linear network. But there is an \emph{implicit assumption} the instructor now makes explicit: the target lies in the span of the eigenfunctions of $T$. If the architecture can only learn linear functions (a deep \emph{linear} network) then what is learned is the best linear approximation, and the loss does not go to zero; if the tangent kernel has the target in the range of its eigenfunctions, we expect learning.
\emph{By choosing the two exponents I can get any error I like. Is that realistic?} The reply: ``we will get to that''; it is not entirely clear how realistic the assumption is, although it is seen in some cases, and what happens with large language models in the feature-learning regime is puzzling.
\emph{Can the same be done for the number of parameters?} No: this is a scaling law in \emph{training time}, which stands in for the number of data points (each step takes fresh data). The parameter count needs a different argument; one can at least count operations (two for the forward pass and four for backpropagation per parameter, giving $C\approx6ND$). ``I am not sure we have a satisfactory answer for the scaling law in the number of parameters''; the heuristic of Section~\ref{sec:L3-params} is the nearest the lecture comes.
\end{discussion}

%% ------------------------------------------------------------------------------------
\section{Does the spectral picture match data?}\label{sec:L3-data}
\tsm{34:41}
For language models one cannot yet compute the necessary spectra. For image classification one can. The instructor describes a study on the CIFAR dataset by Bordelon, Atanasov and Pehlevan\flag\,\citep{bordelon2025}.

\begin{enumerate}
  \item \textbf{Measure $\alpha$.} Compute the eigenvalues of the NTK \emph{at initialization} on the images; they decay like a power law (the finite number of images makes the matrix finite rank, whence a drop at the end of the plot).
  \item \textbf{Measure $s$.} Project the true classification function (cat versus table versus dog, a simple classification task) onto the first $k$ eigenfunctions; the loss \emph{not} explained by the first $k$ modes decays like $k^{-s}$.
  \item \textbf{Predict.} Combine the two to predict the initial slope $t^{-s/\alpha}$ of the training loss of a convolutional network trained on the task: the initial slope matches the observed decay well.
  \item \textbf{What happens next.} Later the NTK starts to \emph{change} (the feature-learning regime). The authors have a theory, \emph{dynamical mean-field theory}, that predicts a second, steeper slope. Finally there is a third stage in which the network learns even \emph{faster}; for this one has no good explanation.
\end{enumerate}
A student asks what the curves labeled by $\gamma$ are; the instructor does not remember.

\begin{beyond}[What that paper claims]
\citet{bordelon2025} build a solvable model that goes beyond the kernel limit and identify three regimes of task difficulty (hard, easy, super-easy). For easy and super-easy targets the exponents are unchanged by feature learning; for \emph{hard} targets feature learning improves the scaling with training time and compute, nearly doubling the exponent. The parameter $\gamma$ in their plots is, as I understand it, the strength of feature learning. Experiments there are on power-law targets on the circle and on CNNs for vision tasks.
\end{beyond}

\paragraph{Language models.}\tsm{41:14}
For large language models it is much less clear what happens with the eigenvalues and with the coefficients $a_k$; the instructor does not think it well understood. He shows a recent paper (from spring or summer)\flag{} that computes the spectrum of the NTK \emph{along training}, not at initialization, and tries to locate the index $k_*$ at which one passes from learned to unlearned eigenfunctions,
measured as the average of the normalized eigenvalues seen by the loss gradient (project the gradient onto the first $k$ eigendirections and see how much remains in the others). It finds a resemblance between how the training loss decreases and how $k_*$ moves, and the smallest learned eigenvalue is smaller in bigger models. Curves seem to rise again at the end, which may be an artifact of the measurement.
\begin{openq}
Nothing like a spectral theory exists that \emph{predicts} the scaling loss of language models. Do the spectra and coefficients of LLMs follow power laws? How do feature learning and the third, faster regime change the exponents?
\end{openq}

%% ------------------------------------------------------------------------------------
\section{Scaling with parameters: the dimension of the data manifold}\label{sec:L3-params}
\tsm{45:20}
The questioner's case: why would the loss fall with the number of parameters $N$? The instructor offers ``a very simple observation'' that relates $N$ for a ReLU network to the \emph{dimension of the data}. Suppose the data are samples from a $d$-dimensional manifold $\mathcal M$, much smaller than the ambient data space (one suspects some structure).

\paragraph{The argument.}\tsm{46:21}
Subdivide $\mathcal M$ into small cells of diameter about $h$. Since $\dim\mathcal M=d$, the number of cells needed is
\[
   M\ \asymp\ h^{-d}\qquad(\text{for a square, }d=2:\ M\,h^2\approx\text{area},\ h\approx M^{-1/2}).
\]
Now take a target $f:\mathcal M\to\R$ that is regular ($C^2$). If one knows the value of $f$ at the cell vertices (and its slope), one approximates it by a \emph{piecewise linear} function, for instance the continuous interpolant, with error of order $h^2$ (Figure~\ref{fig:L3-pl}); the squared error, which is what the loss measures, is of order $h^4$.
A ReLU network defines a piecewise linear function, and to specify a function with different slopes on $M$ different patches one needs, up to a constant depending on the dimension, about $M$ parameters. So $N\approx M\approx h^{-d}$ (up to constants, and for $d\ge2$ up to logarithmic factors), i.e.\ $h\approx N^{-1/d}$, and
\begin{equation}\label{eq:N4d}
   \boxed{\;L(N)\;\lesssim\;h^{4}\;\approx\;N^{-4/d}\;}
\end{equation}
This heuristic relates the intrinsic dimension of the sampled manifold to the decay of the loss in a situation with plenty of data and enough parameters\,\citep{sharma2022}; the lecture takes the exponent $4/d$ as the prediction.

\begin{figure}[tb]
\centering
\begin{tikzpicture}[x=1.88cm, y=1.5cm]
  \draw[ink!45, ->] (-0.15,0) -- (6.35,0) node[below left, font=\sffamily\scriptsize, text=ink] {$x$};
  \foreach \a in {0,1.5,3,4.5} {
    \pgfmathsetmacro{\b}{\a+1.5}
    \pgfmathsetmacro{\fa}{1.0+0.75*sin(deg(\a))+0.12*\a}
    \pgfmathsetmacro{\fb}{1.0+0.75*sin(deg(\b))+0.12*\b}
    \fill[amber!35] (\a,\fa) -- (\b,\fb) -- plot[domain=\b:\a, samples=24, smooth] (\x,{1.0+0.75*sin(deg(\x))+0.12*\x}) -- cycle;
    \draw[slate!50, line width=0.4pt, dotted] (\b,0) -- (\b,\fb);
  }
  \draw[accent, line width=1.8pt, domain=0:6, samples=120, smooth] plot (\x, {1.0+0.75*sin(deg(\x))+0.12*\x});
  \foreach \a in {0,1.5,3,4.5} {
    \pgfmathsetmacro{\b}{\a+1.5}
    \draw[amber!85!black, line width=1.5pt] (\a,{1.0+0.75*sin(deg(\a))+0.12*\a}) -- (\b,{1.0+0.75*sin(deg(\b))+0.12*\b});
    \fill[amber!85!black] (\a,{1.0+0.75*sin(deg(\a))+0.12*\a}) circle (1.1pt);
  }
  \fill[amber!85!black] (6.0,{1.0+0.75*sin(deg(6.0))+0.12*6.0}) circle (1.1pt);
  \draw[<->, ink, line width=0.6pt] (0,-0.17) -- (1.5,-0.17) node[midway, below, font=\sffamily\scriptsize] {cell of size $h$};
  \node[font=\sffamily\scriptsize, text=accent, anchor=south west] at (0.25,1.62) {target $f$};
  \node[font=\sffamily\scriptsize, text=amber!80!black, anchor=north west, align=left] at (3.35,0.55) {shaded: error of the piecewise-linear fit,\\$\lesssim h^{2}$ (squared: $\lesssim h^{4}$)};
\end{tikzpicture}
\caption{Approximating a smooth function by a piecewise-linear one on cells of size $h$ (here $d=1$). A ReLU network with about one parameter per cell can represent the fit, which gives $N\approx h^{-d}$ and $L\lesssim N^{-4/d}$.}
\label{fig:L3-pl}
\end{figure}

\begin{beyond}[The error bound, and what $4/d$ would say]
For $f\in C^2$ on an interval $[a,a+h]$ the linear interpolant $\ell$ through the endpoints satisfies $\abs{f-\ell}\le\tfrac18h^2\max\abs{f''}$ (a standard estimate; the same order holds on cells in higher dimension), so the squared $L^2$ error on $\mathcal M$ is $\lesssim h^4$ and the loss follows \eqref{eq:N4d}. The exponent $2/d$ for the uniform error is the best possible worst-case rate over the unit ball of $C^2$ for approximation by $N$ parameters that depend continuously on $f$ \citep{devore1989}; very deep ReLU networks, whose weights may depend discontinuously on $f$, can do better (up to logarithmic factors, $N^{-4/d}$ in the uniform error\,\citep{yarotsky2018}), and the count $N\approx M$ holds only up to logarithms for $d\ge2$\,\citep{yarotsky2017}. So $4/d$ is the prediction of a specific (shallow, continuous) approximation scheme that matches the observed exponents, not a ceiling on what networks can achieve.
Reading \eqref{eq:N4d} backwards, an empirical exponent $\alpha_N\approx0.076$ (Kaplan et al.) would correspond to an intrinsic dimension $d\approx4/0.076\approx53$; \citet{sharma2022} test the relation for both images and language by estimating $d$ independently.
\end{beyond}

\begin{caution}
Whether language should be viewed as samples from a finite-dimensional manifold is not clear (``perhaps it makes more sense for images''). The argument is for a particular regime, with plenty of data and \emph{enough} parameters. And the empirical laws of Kaplan and Chinchilla concern optimally trained models of increasing size: during training the shapes of the curves change a lot, but once one reaches the optimum, one sees a scaling law across the models.
\end{caution}

% =====================================================================================
%  Lecture 3, part B  —  sections 3.5 – 3.6
% =====================================================================================

%% ------------------------------------------------------------------------------------
\section{Generalization: the classical picture and why it fails}\label{sec:L3-classical}
\tsm{52:48}
Now to the second theme. The question one wants answered: is the network learning \emph{some} function that happens to fit the training points, or \emph{the} function that will work on points it has not seen? Suppose we have a collection of data points (the training set) and then
test how well the fitted model does on other points. The model needs enough \emph{degrees of freedom} to adapt to the data; if it has too few, it may be a poor approximation (\term{underfitting}). If it has very many, it may be very good at matching these particular points and still be a worse approximation than a model with fewer degrees of freedom,
which chooses a simpler fit and ignores part of the noise (\term{overfitting}). So there are two worries: too little complexity in the class of functions, and too much.

\subsection{The bias--variance decomposition}\label{sec:L3-biasvar}\tsm{56:18}
A good way to write this is to decompose the loss. The tool is an elementary identity: for a random variable $Z$ and a constant $c$,
\begin{equation}\label{eq:bv-basic}
  \E\big[(Z-c)^2\big]\;=\;\Var(Z)+\big(\E Z-c\big)^2 ,
\end{equation}
which is proved by adding and subtracting $\E Z$ inside the square: the cross term is $2\,\E[Z-\E Z]\,(\E Z-c)=0$.

\begin{lem}[Bias--variance--noise decomposition]\label{lem:biasvar}
Let $Y=f(x)+\varepsilon$ be a fresh test observation at the fixed input $x$, where $\E\varepsilon=0$, $\Var\varepsilon=\sigma^2$ and $\varepsilon$ is independent of the training sample $S$. Let $\hat f_S$ be a predictor trained on $S$, with $\E_S\hat f_S(x)^2<\infty$. Then
\[
  \E_{S,\varepsilon}\Big[\big(Y-\hat f_S(x)\big)^2\Big]\;=\;\big(\E_S\hat f_S(x)-f(x)\big)^2\;+\;\Var_S\big(\hat f_S(x)\big)\;+\;\sigma^2 .
\]
The three terms are the \emph{squared bias}, the \emph{variance} and the irreducible \emph{noise}.
\end{lem}
\begin{proof}
Fix $S$ and apply \eqref{eq:bv-basic} to $Z=Y$ and $c=\hat f_S(x)$: $\E_\varepsilon[(Y-\hat f_S(x))^2]=\sigma^2+(f(x)-\hat f_S(x))^2$. Now average over $S$ and apply \eqref{eq:bv-basic} to $Z=\hat f_S(x)$ and $c=f(x)$:
$\E_S[(\hat f_S(x)-f(x))^2]=\Var_S(\hat f_S(x))+(\E_S\hat f_S(x)-f(x))^2$.
\end{proof}

The first term says how far the \emph{average} prediction is from the true value, the second how much the prediction \emph{fluctuates} from one sample to the next, the last is noise. The instructor shows the picture from Scott Fortmann-Roe's lecture notes\,\citep{fortmannroe2012} (Figure~\ref{fig:L3-bullseye}): low bias and low variance means hitting the bull's-eye repeatedly; low bias and high variance means the average is right but the shots are scattered;
high bias and low variance is a tight cluster off target; high bias and high variance is scattered and off.

\begin{figure}[tb]
\centering
\begin{tikzpicture}[x=1mm, y=1mm, font=\sffamily\scriptsize]
  \newcommand{\bseye}[5]{% center x, center y, offset x, offset y, spread
    \foreach \r/\c in {15/ink!6, 11.5/ink!11, 8/ink!18, 4.3/crimson!45} \fill[\c] (#1,#2) circle (\r);
    \foreach \dx/\dy in {0.0/0.1, 0.9/0.3, -0.7/0.8, 0.4/-0.95, -0.95/-0.4, 0.6/0.95, -0.2/-0.25, 0.95/-0.45} {
      \fill[accent!85!black] ({#1+#3+#5*\dx},{#2+#4+#5*\dy}) circle (0.85);
    }
  }
  \bseye{24}{22}{0}{0}{2.2}   \bseye{80}{22}{0}{0}{9}
  \bseye{24}{-16}{7}{-4.5}{2.2}   \bseye{80}{-16}{6}{-4}{9}
  \node[anchor=south, font=\headsf\bfseries\scriptsize] at (24,38.5) {low variance};
  \node[anchor=south, font=\headsf\bfseries\scriptsize] at (80,38.5) {high variance};
  \node[anchor=east, font=\headsf\bfseries\scriptsize, rotate=90] at (2.5,22) {low bias};
  \node[anchor=east, font=\headsf\bfseries\scriptsize, rotate=90] at (2.5,-16) {high bias};
\end{tikzpicture}
\caption{Bias and variance (after Fortmann-Roe). Each dot is a prediction from a different training sample; the target is the bull's-eye. Bias is the offset of the cloud's center, variance its spread.}
\label{fig:L3-bullseye}
\end{figure}

\subsection{The classical heuristic}\label{sec:L3-heuristic}\tsm{1:00:37}
In many situations there is a \emph{trade-off} between the two. As one increases the complexity of the class of functions, a minimization over a very complex class can find a function that fits the data points almost exactly, so \emph{bias tends to zero}. But once the complexity is very large the functions are complex enough to hit every training point, which makes them more sensitive to the sample: the solution is more scattered, and the \emph{variance} term grows (Figure~\ref{fig:L3-ucurve}).
The recommended practice is to find the sweet spot.

\begin{figure}[tb]
\centering
\begin{tikzpicture}
\begin{axis}[notes, width=0.78\linewidth, height=0.46\linewidth, xmin=0, xmax=10, ymin=0, ymax=3.4, xlabel={complexity of the function class}, ylabel={expected error}, xtick=\empty, ytick=\empty,
   legend style={at={(0.5,1.02)}, anchor=south, legend columns=4, draw=none, fill=none}]
  \addplot[accent, line width=1.5pt, domain=0:10, samples=80] {2.4*exp(-0.45*x)};
  \addplot[crimson, line width=1.5pt, domain=0:10, samples=80] {0.018*x^2};
  \addplot[ink, line width=2.2pt, domain=0:10, samples=80] {2.4*exp(-0.45*x)+0.018*x^2+0.3};
  \addplot[slate, dashed, line width=1.0pt, domain=0:10, samples=2] {0.3};
  \legend{bias$^2$, variance, total, noise}
  \draw[amber, dashed, line width=0.8pt] (axis cs:4.3,0) -- (axis cs:4.3,1.45);
  \node[font=\sffamily\scriptsize, text=amber!80!black, anchor=south] at (axis cs:4.3,1.5) {sweet spot};
\end{axis}
\end{tikzpicture}
\caption{The classical trade-off (schematic): bias$^2$ falls with complexity, variance rises, and the total error has a minimum at an intermediate complexity.}
\label{fig:L3-ucurve}
\end{figure}

\begin{discussion}[A heuristic, not a law]
A student asks whether this is a heuristic or a theoretical statement. It is a \emph{heuristic that can be proved as a theorem for particular classes} of learning procedures. It governed much of statistical learning theory in the 1990s and 2000s and made people care a lot about limiting the complexity of the class of functions one is allowed to use.
It sits uneasily with scaling: as models get bigger, their complexity increases, so one would suspect that they will not generalize.
There is also a worry about \emph{general intelligence}: it is supposed to solve complicated problems, so it must live in a complicated class of functions, and then the heuristic predicts that it will fit its data too tightly and learn complicated solutions instead of simple ones that generalize.
\end{discussion}

\subsection{Zhang et al.: fitting noise}\label{sec:L3-zhang}\tsm{1:03:13}
Neural networks do not obey this heuristic. The instructor's example is the 2017 paper of Zhang and co-authors\,\citep{zhang2017}. They took the CIFAR image-classification problem and several architectures (convolutional networks and fully connected MLPs), first trained them on the usual task, and then \emph{randomized} things.

\begin{exmp}[Training on noise]
\begin{itemize}
  \item \textbf{Random labels.} The labels of the training images are randomly permuted: dogs are no longer dogs, cats no longer cats. Networks of the same size fit this perfectly: \emph{zero training loss}.
  \item \textbf{Random pixels.} Random perturbations of the pixels of the images: again zero training loss for all variants.
\end{itemize}
In other words, the network has enough capacity to learn \emph{anything} about this data. With random labels the test error is obviously large, since there is nothing to predict (the lecture quotes about $90\%$, chance level for ten classes): no information on unseen data can be extracted.
\end{exmp}

\paragraph{The puzzle.} Two kinds of learning seem to occur. With unstructured random data the model is large and expressive enough to \emph{memorize} everything. When it is given real images, it does not choose one of the vast number of solutions that fit the training data perfectly but do nothing useful on the test data;
it chooses \emph{one that works on the test data too}. There are very many algorithms it could have found with zero training loss, the vast majority with very large test loss, and it does not pick them.
This runs against the heuristic about the trade-off between complexity and variance, and shows that the capacity of the class cannot be what explains generalization: the same class can fit noise.

\begin{beyond}[The lesson in modern terms]
\citet{zhang2017} concluded that traditional capacity measures of the hypothesis class (such as VC dimension or Rademacher complexity) cannot explain generalization in deep networks, since the class fits random labels. What is special must be \emph{which} interpolating solution the training procedure selects (an \emph{implicit bias} of the architecture, optimizer and data). The rest of the lecture is a study of \emph{inductive bias as compression}: why simplicity weighting can give generalization even when the class is huge.
\end{beyond}

%% ------------------------------------------------------------------------------------
\section{Learning as compression}\label{sec:L3-compression}
\tsm{1:07:25}
How could one build an algorithm that learns in a way that generalizes to unseen data, in principle? The lecture takes the road of \emph{sequential prediction}: given $x_1,\dots,x_n$ drawn from a true distribution, predict $x_{n+1}$.

\subsection{Bayesian prediction and log-loss}\label{sec:L3-bayes}\tsm{1:07:55}
Let $\mathcal H$ be a class of hypotheses; each $h\in\mathcal H$ determines a probability distribution $p_h$ over sequences, and $p_h(x_{1:n})$ is the probability it gives to the first $n$ symbols. Fix a \term{prior} $\pi(h)$ on $\mathcal H$, a probability distribution ($\sum_h\pi(h)=1$; Solomonoff's prior below is only a sub-probability and will be normalized).
The standard way to predict is Bayes' rule:
\begin{align}
  \pi(h\mid x_{1:n})&\;\propto\;\pi(h)\,p_h(x_{1:n}) &&\text{(posterior)},\label{eq:posterior}\\
  \Pbayes(x_{n+1}\mid x_{1:n})&\;=\;\sum_{h}\pi(h\mid x_{1:n})\,p_h(x_{n+1}\mid x_{1:n}) &&\text{(prediction)}.\label{eq:predict}
\end{align}
Each hypothesis is weighted by how well it predicted the data seen so far, times its prior weight, and the predictions are averaged; one can view this as an iterative process: after seeing $k$ points one has a posterior, and the next point updates it. Equivalently, the prediction is that of the \term{Bayes mixture} $\Pbayes(x_{1:n})=\sum_h\pi(h)\,p_h(x_{1:n})$.

How do we measure performance? By the \term{log-loss}, the negative logarithm of the probability assigned to the next point by Bayes' rule,
\[
  \ell_t\;=\;-\log\Pbayes(x_t\mid x_{<t})\;=\;\log\frac1{\Pbayes(x_t\mid x_{<t})}\qquad(\text{the ``surprise''}),
\]
which is large when the probability is small. By the chain rule of probability the losses add up to the \emph{total} log-probability of the data: $\sum_{t=1}^n\ell_t=-\log\Pbayes(x_{1:n})$. (This is also what language models minimize, Lecture~2.)

\subsection{Why the logarithm? Prediction is compression}\label{sec:L3-kraft}\tsm{1:11:07}
A student asks why the logarithm, or any particular function, is the right measure. The answer lies in a link between probability and \emph{coding}. Suppose we want to describe each of a collection of events in a fixed alphabet (say binary strings), \emph{briefly}: frequent events should get short descriptions and rare events longer ones, so that the total number of letters used is smallest. To avoid ambiguity one uses
\term{prefix-free} codes: no codeword is the beginning of another.

\begin{thm}[Kraft and Shannon]\label{thm:kraft}
\textbf{(Kraft)} Codeword lengths $\ell(x)$ of a binary prefix-free code satisfy $\sum_x2^{-\ell(x)}\le1$; conversely, any lengths satisfying this inequality are the lengths of some prefix-free code. In particular for any distribution $P$ there is a code with $\ell(x)=\lceil-\log_2P(x)\rceil$.
\textbf{(Shannon)} For data distributed according to $Q$, every prefix-free code has $\E_Q[\ell(X)]\ge\E_Q[-\log_2Q(X)]=:\Hent(Q)$, and using the code built from a different distribution $P$ costs $\E_Q[-\log_2P(X)]=\Hent(Q)+\KL(Q\,\|\,P)$ bits on average (ignoring the rounding of lengths to integers).
\end{thm}
\noindent So the log-loss of a predictor is the \emph{length of the code} it induces, and making good predictions and compressing well are the same problem. (Proofs are in the exercises at the end of the lecture; the lecture states both as known facts.)

\begin{exmp}[A four-letter alphabet]
Let $Q=(\tfrac12,\tfrac14,\tfrac18,\tfrac18)$ on $\{a,b,c,d\}$. The code $a\mapsto0$, $b\mapsto10$, $c\mapsto110$, $d\mapsto111$ is prefix-free; its lengths $1,2,3,3$ equal $-\log_2Q$, and $\sum2^{-\ell}=1$. The expected length is $1.75$ bits $=\Hent(Q)$. Coding the same data with the uniform distribution (length $2$ bits for each letter) costs $2$ bits,
$0.25$ more: this is $\KL(Q\,\|\,U)=0.25$.
\end{exmp}

\paragraph{Kullback--Leibler divergence.}\tsm{1:18:57}
For distributions $Q,P$ on the same space,
\[
   \KL(Q\,\|\,P)\;=\;\E_Q\Big[\log\frac{Q(X)}{P(X)}\Big]\;\ge\;0,
\]
with equality if and only if $P=Q$ (almost everywhere). It is the \emph{expected extra code length} when one switches from the code of the truth $Q$ to the code of $P$. It is not symmetric, but behaves like a distance between distributions.

\paragraph{Regret.}\tsm{1:16:19}
Suppose the data really follow some distribution $Q=p_{h^*}$ with $h^*\in\mathcal H$ (the \emph{realizable} case). Even the true distribution would incur an expected loss equal to the Shannon entropy, which cannot be beaten. What we can ask is how much \emph{worse} the Bayesian prediction is: the \term{regret}
\begin{equation}\label{eq:regret}
   R_n\;=\;\E_{Q}\Big[\log\frac{1}{\Pbayes(x_{1:n})}\Big]-\E_{Q}\Big[\log\frac{1}{Q(x_{1:n})}\Big]\;=\;\KL\big(Q^{(n)}\,\|\,\Pbayes^{(n)}\big),
\end{equation}
the extra code length for the first $n$ points ($Q^{(n)}$, $\Pbayes^{(n)}$ denote the laws of $x_{1:n}$). By the chain rule for the divergence, $R_n=\sum_{t=1}^n\E\,\KL\big(Q(\cdot\mid x_{<t})\,\|\,\Pbayes(\cdot\mid x_{<t})\big)$: the regret is the cumulated expected error of the one-step predictions.

\begin{discussion}[Does the prior matter?]
A student notes that the initial hypothesis is not in the notation. It is there, hidden in the Bayes mixture: it is the prior $\pi$. Does it stop mattering as $n$ grows? ``It depends on the class of hypotheses; I will get to that.''
\end{discussion}

\subsection{Solomonoff's prior and the regret bound}\label{sec:L3-solomonoff}\tsm{1:20:30}
Which class and prior to choose? The instructor chooses an enormous class: \emph{all computable probability distributions}. By the heuristic of the previous section a huge class of hypotheses should lead to bad regret. We will see that it does not.
For the prior one uses \term{Kolmogorov complexity}: for a hypothesis $h$, $\Kc(h)$ is the length of the shortest (prefix-free, for technical reasons) program that computes it on a fixed universal machine. Set
\[
   w(h)\;=\;2^{-\Kc(h)}\qquad(\text{Solomonoff's prior}).
\]
By Kraft's inequality $Z:=\sum_hw(h)\le1$, so $w$ is a sub-probability; the prior used in \eqref{eq:posterior}--\eqref{eq:predict} is its normalization $\pi=w/Z$. It is \term{Occam's razor} made quantitative: among the hypotheses that describe the data, simplicity, measured by program length, is rewarded.

\begin{thm}[Regret is bounded by the Kolmogorov complexity of the truth]\label{thm:solomonoff}
Let $\mathcal H$ be the (countable) class of computable distributions with the prior $\pi=w/Z$, $w(h)=2^{-\Kc(h)}$, and suppose the data are drawn from $Q=p_{h^*}$ with $h^*\in\mathcal H$. Then for \emph{every} $n$ and \emph{every} sequence $x_{1:n}$,
\[
   \log_2\frac{Q(x_{1:n})}{\Pbayes(x_{1:n})}\;\le\;\log_2\frac1{\pi(h^*)}\;=\;\Kc(h^*)+\log_2Z\;\le\;\Kc(h^*),
\]
and hence $R_n\le\Kc(h^*)$ bits for all $n$.
\end{thm}
\begin{proof}
All terms in the Bayes mixture are non-negative, so it is at least the term of the true hypothesis: $\Pbayes(x_{1:n})=\sum_h\pi(h)p_h(x_{1:n})\ge\pi(h^*)\,p_{h^*}(x_{1:n})=\frac{2^{-\Kc(h^*)}}{Z}\,Q(x_{1:n})\ge2^{-\Kc(h^*)}\,Q(x_{1:n})$, because $Z\le1$. Take logarithms.
\end{proof}
\noindent In the words of the lecture: after $n$ observations the worst that can happen is the Shannon entropy, which cannot be reduced, \emph{plus} one extra term, the Kolmogorov complexity of the true hypothesis. If what we are trying to learn is a computable function with a short description, then this argument ``will certainly get us there''.

\begin{intuition}
The bound does not grow with $n$, and apart from the prior weight $\pi(h^*)$ of the truth it does not depend on how large or complex the class is (for a finite class with the uniform prior it is $\log_2\abs{\mathcal H}$). The prior does the work: it spends its probability mass mostly on simple hypotheses, so that a simple truth is quickly recognized, whereas the (infinitely many) complicated hypotheses can wait in the tail. Compare with the classical heuristic, in which complexity of the class itself is penalized; here the class is unrestricted and the penalty is \emph{soft}.
\end{intuition}

\paragraph{Consequences.} By the chain rule, $R_n=\sum_{t\le n}\E\,\KL_t\le\Kc(h^*)$ for all $n$, where $\KL_t=\KL\big(Q(\cdot\mid x_{<t})\,\|\,\Pbayes(\cdot\mid x_{<t})\big)$ is the one-step error. So the series $\sum_t\E\,\KL_t$ converges: the expected number of steps with $\KL_t>\epsilon$ is at most $\Kc(h^*)/\epsilon$ (the Bayes predictor can make large errors only finitely often, in expectation; the terms themselves need not be monotone), and the average regret $R_n/n\to0$. A student asks whether this is ``less than that for all $n$'': yes, the regret after $n$ data points is bounded by the same constant, so the per-point regret vanishes, ``in a sense, I am learning''.

\begin{caution}
The difficulty is that \emph{$\Kc$ is not computable}. The hope is that, for the kinds of problems neural networks learn, short programs exist, and that what a neural network does is effectively \emph{to give weight to solutions in a way that penalizes longer ones}. That is regularization: ``in your reading material there is a link to a paper that calls this the \emph{soft inductive bias}'', a penalty for more complicated programs closely related to penalizing large weights in the network.
\end{caution}

\begin{beyond}[Two literature pointers]
The argument is Solomonoff's\,\citep{solomonoff1964,solomonoff1978} (his actual construction mixes over the class of \emph{enumerable semimeasures}, which makes the mixture itself enumerable; the proof above uses only the domination $\Pbayes\ge\pi(h^*)p_{h^*}$ and goes through verbatim); see \citet{li2019} for a complete treatment of Kolmogorov complexity, and \citet{grunwald2007} for the closely related MDL principle. The reading that calls it a \emph{soft inductive bias} is, I believe, \citet{wilson2025}, who argues that benign overfitting, double descent and the generalization of large models are not as mysterious as often thought, if inductive biases are soft preferences rather than hard restrictions.
\end{beyond}

% =====================================================================================
%  Lecture 3, part C  —  section 3.7 and closing material
% =====================================================================================

%% ------------------------------------------------------------------------------------
\section{Continuous parameters: from Rissanen to singular learning theory}\label{sec:L3-singular}
\tsm{1:32:26}
What happens in a situation closer to neural networks: a large parameter space $\Theta\subset\R^d$, with each $\theta$ giving a distribution $p_\theta$? Instead of a countable collection of hypotheses we now have a \emph{continuum}, and the Bayesian story is the same with integrals in place of sums.

\subsection{Evidence and regret}\label{sec:L3-evidence}\tsm{1:33:27}
Assume that the truth lies in the family: $Q=p_{\theta^*}$ for some $\theta^*\in\Theta$. Start from a prior density $\pi(\theta)$ and, given $x_1,\dots,x_n$, form the posterior
\[
   \pi_n(\theta)\;=\;\frac{\pi(\theta)\prod_{i=1}^np_\theta(x_i)}{Z_n},\qquad Z_n=\int\pi(\theta)\prod_{i=1}^np_\theta(x_i)\,\dd\theta ,
\]
where the \emph{normalizing constant} $Z_n$ (the \term{evidence} or marginal likelihood) turns out to be an important object. It says how \emph{surprised} the model, with its prior over parameters, should be by the data: it is the probability the model assigns to the data. And, exactly as before,
$-\log Z_n=\sum_{t=1}^n-\log\Pbayes(x_t\mid x_{<t})$ is the code length of the data under the Bayesian predictions, and the regret is
\[
   R_n\;=\;\E_{p_{\theta^*}}\Big[\log\frac{p_{\theta^*}(x_{1:n})}{Z_n}\Big].
\]

\begin{thm}[Rissanen; Clarke--Barron]\label{thm:rissanen}
Let $\{p_\theta\}_{\theta\in\Theta\subset\R^d}$ be a regular parametric family: identifiable and smooth, with $\theta^*$ an interior point where the Fisher information matrix $I(\theta^*)=\nabla^2_\theta\KL(p_{\theta^*}\|p_\theta)\big|_{\theta=\theta^*}$ is positive definite, and let the prior density be continuous and positive at $\theta^*$; assume also the technical regularity conditions of \citet{clarke1990} (smoothness and moment bounds for $\log p_\theta$ near $\theta^*$, a proper prior, and consistency of the posterior, for instance for compact $\Theta$). For i.i.d.\ data from $p_{\theta^*}$,
\[
  R_n\;=\;\frac d2\,\log\frac{n}{2\pi e}\;+\;\log\frac{\sqrt{\det I(\theta^*)}}{\pi(\theta^*)}\;+\;o(1)\qquad(n\to\infty).
\]
\end{thm}
\noindent This is the content of the result the lecture credits to Rissanen\,\citep{rissanen1986,rissanen1996}; the precise form with these constants is due to \citet{clarke1990}. As $n\to\infty$ the posterior looks more and more like a Gaussian centered at the true parameter and converges to a delta function there; the regret grows like $\tfrac d2\log n$, \emph{depending on the dimension $d$ of the parameter space}.
The assumptions are essential: the truth is in the model, and the divergence $\KL(p_{\theta^*}\|p_\theta)$ has a \emph{positive definite} Hessian at $\theta^*$.

\begin{intuition}
Each parameter must be specified to a precision of order $n^{-1/2}$ (the statistical resolution of $n$ observations); that costs about $\tfrac12\log n$ nats per coordinate, which is the origin of the BIC/MDL penalty $\tfrac d2\log n$. Compare with Theorem~\ref{thm:solomonoff}: for a \emph{countable} class with a simplicity prior the regret is a \emph{constant}; for a \emph{continuous} parametric family it grows with $n$ and with $d$.
\end{intuition}

The instructor's verdict is that this is ``bad news compared with Solomonoff induction'': the regret increases with the complexity (dimension) of the parameter space. Two caveats keep the news from being so bad for neural networks.

\begin{enumerate}
  \item \textbf{Not in the regime.} The theorem is a statement about $n\gg d$. For language models the amount of data is typically larger than the number of parameters, but probably still not by enough to be in the asymptotic regime where this behavior appears; in image experiments such as that of Zhang et al., the networks have more parameters than training examples.
  \item \textbf{The Hessian is not positive definite.} Neural networks have many directions along which the loss does not change, so the assumption of a positive definite Hessian fails, at least at realistic error levels. And there is typically not one true parameter but a whole \emph{set} of parameters that give the true distribution.
\end{enumerate}

\begin{figure}[tb]
\centering
\begin{tikzpicture}
\begin{axis}[notes, width=0.88\linewidth, height=0.56\linewidth, xmode=log, xmin=1, xmax=1e5, ymin=0, ymax=15, xlabel={number of observations $n$}, ylabel={regret $R_n$ (bits)},
   legend style={at={(0.03,0.97)}, anchor=north west, font=\sffamily\fontsize{6.2}{7}\selectfont}]
  \addplot[crimson, line width=1.8pt] table[x=n, y=Rfinite] {figures/data/regret.dat};   \addlegendentry{finite class of 16, uniform prior (exact)}
  \addplot[accent, line width=1.8pt] table[x=n, y=Rcont] {figures/data/regret.dat};   \addlegendentry{Bernoulli, uniform prior (exact)}
  \addplot[ink, dashed, line width=1.0pt] table[x=n, y=Rasym] {figures/data/regret.dat};   \addlegendentry{$\frac12\log_2\frac n{2\pi e}+\frac12\log_2\frac1{\theta^*(1-\theta^*)}$}
  \draw[crimson, dotted, line width=0.9pt] (axis cs:1,4) -- (axis cs:1e5,4);
  \node[font=\sffamily\scriptsize, text=crimson, anchor=south east] at (axis cs:9e4,4.1) {$\log_2 16=4$};
\end{axis}
\end{tikzpicture}
\caption{Regret of Bayesian prediction in two models of coin flips with true bias $\theta^*=5/17$ (exact computation by summing over the binomial). \textbf{Finite class} $\{j/17\}_{j=1}^{16}$ with uniform prior: the regret saturates at $\log_2 16=4$ bits, the bound of Theorem~\ref{thm:solomonoff} (with $\Kc$ replaced by $-\log_2\pi$). \textbf{One continuous parameter}, uniform prior: it keeps growing like $\frac12\log_2n$ and tracks the formula of Theorem~\ref{thm:rissanen} (dashed); at $n=10^5$ the two differ by $1.5\cdot10^{-5}$ bits.}
\label{fig:L3-regret}
\end{figure}

\subsection{Singular learning theory}\label{sec:L3-watanabe}\tsm{1:41:55}
There is a different result of Watanabe\,\citep{watanabe2009,watanabe2018} which does \emph{not} assume that the Hessian is positive definite. Let $K(\theta)=\KL(Q\,\|\,p_\theta)\ge0$ and let $W_0=\{\theta:K(\theta)=0\}$ be the set of true parameters; it can be a complicated (algebraic or analytic) subset, not a single point.
When the Hessian has zero eigenvalues, the regret depends in an interesting way on \emph{how quickly $K$ grows as one moves away from $W_0$}, and the leading term is governed by a number called the \term{real log canonical threshold} (RLCT).

\begin{defn}[Volume scaling and the RLCT]\label{def:rlct}
Let $V(\epsilon)=\int_{\{K(\theta)<\epsilon\}}\pi(\theta)\,\dd\theta$ be the prior mass of the parameters whose divergence is below $\epsilon$. For analytic $K$ there are $\rlct>0$ and an integer $m\ge1$ with
\[
   V(\epsilon)\;\sim\;c\;\epsilon^{\,\rlct}\,\big(\log\tfrac1\epsilon\big)^{m-1}\qquad(\epsilon\to0).
\]
$\rlct$ is the \emph{real log canonical threshold} and $m$ its \emph{multiplicity}. It measures \emph{flatness} around $W_0$: how the volume below level $\epsilon$ scales as one moves away from the truth.
\end{defn}

\begin{thm}[Watanabe]\label{thm:watanabe}
In the realizable case, under Watanabe's regularity conditions (a compact parameter set, a positive analytic prior, a log-likelihood ratio that is analytic in $\theta$ with values in $L^s(Q)$, and relatively finite variance, $\E_Q[(\log\frac{q}{p_\theta})^2]\le c\,K(\theta)$),
\[
   R_n\;=\;\rlct\,\log n\;-\;(m-1)\log\log n\;+\;O(1).
\]
For a regular model $\rlct=d/2$ and $m=1$, which recovers the leading term of Theorem~\ref{thm:rissanen}; in general $\rlct\le d/2$.
\end{thm}
\noindent (The lecture states only the leading behavior: ``the principal order term depends on the real log canonical threshold''.)

\paragraph{The simplest case.}\tsm{1:44:44}
Take the case in which $W_0$ is a smooth submanifold of dimension $d-r$ and $K$ is a \emph{Morse--Bott} function: zero on $W_0$ and with positive definite Hessian in the $r$ directions normal to it. Then near $W_0$, $K\approx\abs{y}^2$ in normal coordinates $y\in\R^r$, so $V(\epsilon)\asymp\epsilon^{r/2}$ and $\rlct=r/2$. The \emph{effective dimension} is the \emph{codimension} of the set of true parameters, and the regret is
$\tfrac r2\log n$ rather than $\tfrac d2\log n$: the $d-r$ parameters that are free to move along $W_0$ cost nothing. The instructor sees in this ``another reason why in some situations neural networks learn surprisingly well''.

\begin{figure}[tb]
\centering
\begin{tikzpicture}
\begin{groupplot}[group style={group size=3 by 1, horizontal sep=3.5mm}, width=0.355\linewidth, height=0.355\linewidth, view={0}{90},
    colormap={kl}{rgb255=(20,33,61) rgb255=(60,98,180) rgb255=(238,242,250)}, point meta min=0, point meta max=1,
    axis lines=box, ticks=none, xmin=-1, xmax=1, ymin=-1, ymax=1, title style={font=\sffamily\scriptsize\bfseries, text=ink, yshift=-1mm}]
\nextgroupplot[title={(a) regular}]
  \addplot3[surf, shader=interp, samples=34, domain=-1:1, y domain=-1:1] {x^2+y^2};
  \fill[amber] (axis cs:0,0) circle (2.2pt);
\nextgroupplot[title={(b) Morse--Bott}]
  \addplot3[surf, shader=interp, samples=34, domain=-1:1, y domain=-1:1] {y^2};
  \draw[amber, line width=1.8pt] (axis cs:-1,0) -- (axis cs:1,0);
\nextgroupplot[title={(c) singular}]
  \addplot3[surf, shader=interp, samples=34, domain=-1:1, y domain=-1:1] {x^2*y^2};
  \draw[amber, line width=1.8pt] (axis cs:-1,0) -- (axis cs:1,0);
  \draw[amber, line width=1.8pt] (axis cs:0,-1) -- (axis cs:0,1);
\end{groupplot}
\end{tikzpicture}
\caption{Three divergence landscapes $K(\theta)$ on $[-1,1]^2$ (dark $=$ small; saturated at $K=1$); the set $W_0$ of true parameters is drawn in amber. \textbf{(a)} $K=\theta_1^2+\theta_2^2$: an isolated minimum, $\rlct=d/2=1$. \textbf{(b)} $K=\theta_2^2$: a whole valley; $\rlct=\tfrac12$ (the codimension over two). \textbf{(c)} $K=\theta_1^2\theta_2^2$: a cross of true parameters with a singular crossing; $\rlct=\tfrac12$ and $m=2$.}
\label{fig:L3-landscapes}
\end{figure}

\begin{beyond}[Worked RLCT examples, and why they matter]
Volumes (Lebesgue areas) on $[-1,1]^2$ for $\epsilon\to0$; a uniform probability prior divides them by $4$ and changes neither $\rlct$ nor $m$:
\begin{itemize}
  \item $K=\theta_2^2$: $V(\epsilon)=4\sqrt\epsilon$, so $\rlct=\tfrac12$, $m=1$ (versus $d/2=1$).
  \item $K=\theta_1^2\theta_2^2$: $\{|\theta_1\theta_2|<\sqrt\epsilon\}$ has area $4\sqrt\epsilon\,(1+\tfrac12\ln\tfrac1\epsilon)$, so $\rlct=\tfrac12$, $m=2$.
  \item $K=\theta_1^2+\theta_2^4$: the region is an ellipse-like set of area $\asymp\epsilon^{1/2}\epsilon^{1/4}$, so $\rlct=\tfrac34$, $m=1$.
\end{itemize}
Watanabe also proves, as a separate theorem, that the Bayes generalization error satisfies $\E\,\KL(Q\,\|\,\Pbayes(\cdot\mid x_{1:n}))=\rlct/n+o(1/n)$ (the regret is the sum of these errors over the first $n$ steps), so \emph{$\rlct$ plays the role of $d/2$} in the familiar $d/(2n)$ of regular models. Over-parameterized networks have large degenerate sets $W_0$ (continuous rescaling symmetries, removable or dead units), so $\rlct$ can be much smaller than $d/2$; merely permuting hidden units only produces finitely many copies of $W_0$ and changes neither $\rlct$ nor $m$.
Estimating a ``local learning coefficient'' for real networks is an active research direction\,\citep{lau2023}. The caveat is that all of this concerns the \emph{Bayesian posterior}, not the output of SGD.
\end{beyond}

\begin{table}[tb]
\centering\small
\renewcommand{\arraystretch}{1.3}
\begin{tabular}{@{}>{\raggedright\arraybackslash}p{33mm}>{\raggedright\arraybackslash}p{23mm}>{\raggedright\arraybackslash}p{39mm}>{\raggedright\arraybackslash}p{16mm}@{}}
\toprule
\textbf{Hypothesis class} & \textbf{Prior} & \textbf{Regret $R_n$} & \textbf{Credit}\\
\midrule
Countable, e.g.\ all computable distributions & $2^{-\Kc(h)}$ & $\le\Kc(h^*)$ bits, for all $n$ & Solomonoff\\
Regular parametric, $\Theta\subset\R^d$ & smooth density & $\tfrac d2\log\frac n{2\pi e}+O(1)$ & Rissanen; Clarke--Barron\\
Singular parametric & smooth density & $\rlct\log n-(m-1)\log\log n+O(1)$, with $\rlct\le d/2$ & Watanabe\\
\bottomrule
\end{tabular}
\caption{Three regimes of Bayesian regret met in this lecture.}
\end{table}

%% ------------------------------------------------------------------------------------
\tocsection{Where this leaves us}
The picture assembled in these three lectures is partial but coherent. Capabilities follow smooth trends (Lecture~1); those trends can be derived, in a solvable model, from the spectrum of a kernel (Lecture~2); and large, apparently over-parameterized models can nevertheless generalize
because learning is a form of \emph{compression}, in which simplicity is rewarded softly and effective complexity is smaller than parameter count (Lecture~3). The instructor stops here, with the remark that grokking and in-context learning, the other items in the opening table, wait for later.

\begin{openq}[Questions this lecture leaves open]
\begin{itemize}
  \item Does SGD on a neural network implement something like a simplicity prior $2^{-\Kc}$? What is the right notion of ``program length'' for a network?
  \item Can the RLCT (or a practical proxy) be measured for large language models, and does it predict generalization or capability jumps?
  \item What is the spectral theory of scaling for language models: do the NTK spectra and target coefficients follow power laws, and what changes in the feature-learning regime, including the ``third'', faster stage?
  \item For safety: can one connect the abrupt transitions (grokking, in-context learning) to a notion like a change in effective dimension, and hence \emph{forecast} sudden capability jumps?
\end{itemize}
\end{openq}

\begin{takeaways}
\begin{itemize}
  \item Open puzzles of deep learning: scaling, generalization, grokking, shared and linear representations, in-context learning. Only for the first two is there a (partial) theory today.
  \item The spectral law $L(t)\asymp t^{-s/\alpha}$ assumes a frozen kernel and a target in the range of the NTK; it is architecture-agnostic but its exponents must be measured. In CIFAR experiments the initial slope is predicted well, later regimes are not.
  \item A heuristic for model size: data on a $d$-dimensional manifold and ReLU networks give $L\propto N^{-4/d}$.
  \item \textbf{Bias--variance:} $\E[(Y-\hat f)^2]=\text{bias}^2+\text{variance}+\sigma^2$. The classical heuristic (complexity $\Rightarrow$ overfitting) fails for deep networks: they fit random labels (zero training loss) yet generalize on real data.
  \item \textbf{Prediction is compression:} log-loss $=$ code length (Kraft, Shannon); regret $=\KL$ between the truth and the Bayes mixture. With Solomonoff's prior, $R_n\le\Kc(h^*)$ for all $n$, however huge the class.
  \item Continuous parametric models pay $\tfrac d2\log n$ (Rissanen); singular models pay only $\rlct\log n$ with $\rlct\le d/2$ (Watanabe), so degenerate parameter space means \emph{lower effective complexity}.
\end{itemize}
\end{takeaways}

\begin{exercise}[Problems]
\begin{enumerate}
  \item \emph{Gibbs' inequality.} Show $\KL(Q\|P)\ge0$ with equality iff $P=Q$. (Hint: $-\log$ is convex; Jensen.)
  \item \emph{Kraft.} Prove $\sum_x2^{-\ell(x)}\le1$ for a binary prefix-free code (hint: the codewords are leaves of a binary tree; count the fraction of the tree below each), and construct a prefix code with lengths $\lceil-\log_2P(x)\rceil$. Deduce $\E_Q[\ell]\ge\Hent(Q)$ using Gibbs.
  \item \emph{Bias--variance for shrinkage.} Observe $y_1,\dots,y_n$ i.i.d.\ with mean $\mu$ and variance $\sigma^2$, and use the estimator $\hat\mu_c=c\,\bar y$ with $c\in[0,1]$. Compute bias$^2$ and variance, the minimizing $c^*$ and the minimal total error. Show that $c^*<1$ whenever $\sigma^2>0$, and that the gain over $c=1$ is largest when $n\mu^2/\sigma^2$ is small.
  \item \emph{Finite classes.} For a finite class with prior $\pi$ and truth $h^*$, show $R_n\le\log\frac1{\pi(h^*)}$ and $\lim_nR_n=\log\frac1{\pi(h^*)}$ (when the hypotheses are distinguishable).
  \item \emph{Laplace's rule.} For coin flips with the uniform prior on $\theta$, show $Z_n=\frac{1}{(n+1)\binom nk}$ for a sequence with $k$ ones, and deduce $R_n$. Check numerically against Theorem~\ref{thm:rissanen} with $d=1$.
  \item \emph{RLCT.} Compute $\rlct$ and $m$ for $K(\theta)=\theta_1^2+\theta_2^4$ and for $K(\theta)=\theta_1^4\theta_2^2$ on $[-1,1]^2$.
  \item \emph{Manifold dimension.} With $d=2$ what is the exponent in \eqref{eq:N4d}? What intrinsic dimension gives an exponent $0.076$?
  \item \emph{Random labels.} Why is the \emph{training} loss of Zhang et al.'s random-label experiment no evidence against the soft-inductive-bias picture? What would the picture predict for the test loss?
\end{enumerate}
\end{exercise}


